% Generated draft: requires mathematical and source review.
\documentclass[11pt]{article}
\usepackage[margin=25mm]{geometry}
\usepackage{fontspec}
\setmainfont{Latin Modern Roman}
\usepackage{amsmath,amssymb,mathtools,unicode-math}
\setmathfont{Latin Modern Math}
\usepackage{enumitem,needspace,xurl,hyperref}
\setlength{\parindent}{0pt}
\setlength{\parskip}{5pt}
\setlength{\emergencystretch}{3em}
\newcommand{\N}{\mathbb{N}}
\newcommand{\Z}{\mathbb{Z}}
\newcommand{\Q}{\mathbb{Q}}
\newcommand{\R}{\mathbb{R}}
\newcommand{\C}{\mathbb{C}}
\newcommand{\F}{\mathbb{F}}
\newcommand{\A}{\mathbb{A}}
\newcommand{\PP}{\mathbb{P}}
\newcommand{\E}{\mathbb{E}}
\newcommand{\eps}{\varepsilon}
\newcommand{\dd}{\,\mathrm{d}}
\begin{document}
\section*{UnsolvedMath ID: 30004935}
\textbf{Source problem code:} catalogue entry
\subsection{Definitions and mathematical statement}
Let \(M\) be a closed smooth manifold, meaning compact and without boundary, and let \(\mathcal R(M)\) denote the ambient space of Riemannian metrics, equipped with the \(C^2\) topology. The supplied record does not specify whether this ambient space consists of smooth metrics or metrics of a fixed finite differentiability class; this choice remains to be fixed, with at least \(C^2\) regularity required here.

For \(g\in\mathcal R(M)\), write
\[
S_gM=\{(x,v)\in TM:g_x(v,v)=1\}
\]
for its unit tangent bundle, and let \(\varphi_g^t\) be the geodesic flow with generator \(X_g\). A closed geodesic is called hyperbolic if the derivative of its local Poincar\'e return map on a section transverse to the corresponding periodic orbit in \(S_gM\) has no complex eigenvalue of modulus \(1\). This transverse definition removes the flow direction. Define
\[
\mathcal R_{\mathrm{hyp}}(M)
=\{g\in\mathcal R(M):\text{every closed geodesic of }g\text{ is hyperbolic}\}.
\]
The flow \(\varphi_g^t\) is Anosov if there is a continuous \(D\varphi_g^t\)-invariant splitting
\[
T(S_gM)=E^s\oplus\R X_g\oplus E^u
\]
and constants \(C\geq1\) and \(\lambda>0\) such that, for every \(z\in S_gM\), every \(t\geq0\), and every \(v_s\in E_z^s\), \(v_u\in E_z^u\),
\[
\|D\varphi_g^t(v_s)\|\leq C e^{-\lambda t}\|v_s\|,
\qquad
\|D\varphi_g^{-t}(v_u)\|\leq C e^{-\lambda t}\|v_u\|,
\]
using any fixed auxiliary Riemannian norm on \(T(S_gM)\).

\textbf{Conjecture as formulated in the supplied registry.} For every closed manifold \(M\), in arbitrary dimension, is
\[
\operatorname{Int}_{C^2}\mathcal R_{\mathrm{hyp}}(M)
=
\{g\in\mathcal R(M):\varphi_g^t\text{ is Anosov}\}?
\]
Here membership in the left-hand side means that some \(C^2\)-open neighborhood of \(g\), relative to \(\mathcal R(M)\), consists entirely of metrics whose closed geodesics are hyperbolic.

\textbf{Unresolved source qualifications.} The original excerpt conjectures the extension of an unspecified ``Theorem A'' to Riemannian metrics on closed manifolds of arbitrary dimensions. That theorem is not supplied, so its equivalence to the displayed registry formulation cannot be established from this record. No lower dimension bound is supplied; in particular, the intended treatment of dimensions \(0\) and \(1\), including whether the stable and unstable bundles must be nonzero, remains unspecified. The displayed splitting imposes no nonzero-rank requirement on those bundles. These conventions require review before an unqualified statement in every dimension is accepted.
\subsection{Short English statement}
The supplied conjecture asks whether a Riemannian metric on a closed manifold has an Anosov geodesic flow exactly when every sufficiently small \(C^2\) perturbation has only hyperbolic closed geodesics. The ambient metric regularity, low-dimensional conventions, and identification with the report's unspecified Theorem A require review.
\subsection{Sources}
\begin{itemize}
\item Supplied local UnsolvedMath registry record, permanent ID \(30004935\), source problem code \texttt{catalogue entry}, title ``Stable Hyperbolicity of Closed Geodesics for Riemannian Metrics,'' in the Oberwolfach Reports Open Problems collection. No individual proposer is named; the supplied proposed year is 2021.
\item Supplied report citation: \emph{Mini-Workshop: \((Anosov)^3\)} (2022), Oberwolfach Reports, pp.\ 3067--3099, DOI \texttt{10.4171/owr/2021/56}. \url{https://doi.org/10.4171/owr/2021/56}. The supplied original excerpt reads: ``We conjecture that Theorem A also holds for Riemannian metrics on closed manifolds of arbitrary dimensions.'' Theorem A itself is absent from the supplied record.
\item Supplied literature source: ``Mini-Workshop: Anosov cubed report,'' \url{https://ems.press/content/serial-article-files/46934}.
\item Literature assessment dated exactly \texttt{2026-08-22}: the supplied curation labels the problem open, reports a corresponding generic statement for Finsler metrics and for Riemannian surfaces, and says that the arbitrary-dimensional closed-Riemannian extension remains conjectural and that no authoritative resolution was found. This is supplied dated curation, not an independently verified literature assessment.
\item The registry labels the statement \texttt{corrected\_verified} and says it was checked against the official report. This is a supplied verification claim; the report has not been independently consulted for this draft.
\end{itemize}
\end{document}
